
\chapter{راست نو و پوپولیسم معاصر}

\begin{tcolorbox}[
    enhanced,
    colback=VeryLightGray, % FIXED: use defined color name
    colframe=RightBlue, % FIXED: use defined color name
    boxrule=2pt,
    arc=8pt,
    fonttitle=\bfseries\large,
    title={\rl{نقل‌قول آغازین}},
    attach boxed title to top center={yshift=-3mm},
    boxed title style={colback=RightBlue, colframe=RightBlue} % FIXED: use defined color name
]
\begin{center}
\Large\textit{«سیاست در پایین‌دست فرهنگ قرار دارد.»}

\vspace{3mm}
\normalsize
--- اندرو برایتبارت (منسوب به آنتونیو گرامشی)
\end{center}
\end{tcolorbox}

\vspace{5mm}

\section*{درآمد}


در دههٔ ۲۰۱۰، جهان سیاست شاهد پدیده‌هایی شد که بسیاری را غافلگیر کرد: پیروزی دونالد ترامپ، رأی برگزیت، صعود احزاب پوپولیست در اروپا. تحلیلگران از «راست نو»، «آلت‌رایت»، «ناسیونالیسم جدید»، و «پوپولیسم» سخن گفتند.

این جریان‌ها چه هستند؟ آیا با راست سنتی تفاوت دارند؟ آیا صرفاً بازگشت فاشیسم‌اند یا پدیده‌ای نوین؟ در این فصل به بررسی این جریان‌های پیچیده و متنوع می‌پردازیم.


%═══════════════════════════════════════════════════════════════
\section{راست نو اروپایی (Nouvelle Droite)}
%═══════════════════════════════════════════════════════════════

\subsection{آلن دوبنوا و GRECE}

\begin{tcolorbox}[
    enhanced,
    colback=gray!5,
    colframe=darkgray,
    boxrule=2pt,
    arc=10pt,
    shadow={2mm}{-2mm}{0mm}{black!30},
    fonttitle=\bfseries\large,
    title={\rl{👤 پرترهٔ متفکر: آلن دوبنوا}}
]

\textbf{آلن دوبنوا} (متولد ۱۹۴۳)، فیلسوف و نویسندهٔ فرانسوی، بنیان‌گذار «راست نو» اروپایی.

\textbf{زندگی:}
\begin{itemize}
    \item در جوانی فعال در گروه‌های ناسیونالیست
    \item تأسیس GRECE در ۱۹۶۸
    \item نویسندهٔ بیش از ۵۰ کتاب
    \item سردبیر مجلات فکری متعدد
\end{itemize}



\textbf{ایده‌های کلیدی:}
\begin{itemize}
    \item متاپولیتیک: تغییر فرهنگ پیش از سیاست
    \item ضد مسیحیت: بازگشت به بت‌پرستی اروپایی
    \item ضد آمریکا و ضد لیبرالیسم
    \item تکثرگرایی قومی و فرهنگی
    \item اروپا به مثابه تمدن
\end{itemize}

\end{tcolorbox}


در ۱۹۶۸، همان سالی که چپ در خیابان‌های پاریس انقلاب می‌کرد، گروهی از روشنفکران راست «گروه تحقیق و مطالعهٔ تمدن اروپایی» (GRECE) را تأسیس کردند. این گروه به «راست نو» (Nouvelle Droite) شهرت یافت.

راست نو از فاشیسم و نازیسم فاصله گرفت. دوبنوا معتقد بود که شکست فاشیسم نشان داد راست نمی‌تواند با خشونت پیروز شود. باید استراتژی جدیدی اتخاذ کرد: \textbf{متاپولیتیک}.

\textbf{متاپولیتیک} (Metapolitics) یعنی تغییر فرهنگ و ذهنیت جامعه پیش از تلاش برای کسب قدرت سیاسی. این ایده از گرامشی مارکسیست وام گرفته شده بود: «هژمونی فرهنگی» پیش‌شرط قدرت سیاسی است.

راست نو به جای فعالیت حزبی، روی نشر کتاب، مجله، و اتاق‌های فکر تمرکز کرد. هدف: تغییر پارادایم فکری جامعه.


\subsection{ایده‌های کلیدی راست نو}


\textbf{ضد مسیحیت:}
دوبنوا معتقد است مسیحیت ریشهٔ همهٔ بدی‌های مدرنیته است: مساوات‌طلبی، جهانشمول‌گرایی، و رد سلسله‌مراتب طبیعی. او به بت‌پرستی اروپایی (یونان و روم و اسکاندیناوی) به عنوان جایگزین نگاه می‌کند.

\textbf{ضد لیبرالیسم:}
لیبرالیسم، چه اقتصادی و چه سیاسی، دشمن اصلی است. لیبرالیسم همهٔ پیوندهای جمعی را منحل می‌کند و جهان را به «بازار جهانی» تبدیل می‌کند.

\textbf{ضد آمریکا:}
آمریکا تجسم لیبرالیسم و مصرف‌گرایی است. اروپا باید از سلطهٔ فرهنگی و سیاسی آمریکا آزاد شود.

\textbf{تفاوت‌گرایی (Différentialisme):}
همهٔ فرهنگ‌ها برابر اما متفاوت‌اند. هر فرهنگی حق دارد در سرزمین خود حفظ شود. این ایده، که گاه «نژادپرستی بدون نژاد» نامیده می‌شود، جایگزین نژادپرستی بیولوژیک شده است.

\textbf{دموکراسی ارگانیک:}
دموکراسی لیبرال که بر فرد استوار است باید جای خود را به دموکراسی مبتنی بر جماعت‌های طبیعی (خانواده، قوم، ملت) بدهد.


%═══════════════════════════════════════════════════════════════
\section{آلت‌رایت آمریکایی}
%═══════════════════════════════════════════════════════════════

\begin{tcolorbox}[
    enhanced,
    colback=rightblue!8,
    colframe=rightblue,
    boxrule=1.5pt,
    arc=5pt,
    fonttitle=\bfseries,
    title={\rl{🔵 تعریف کلیدی: آلت‌رایت}}
]
\textbf{آلت‌رایت} (Alt-Right یا راست جایگزین) جنبشی سیاسی و فرهنگی است که در دههٔ ۲۰۱۰ در آمریکا ظهور کرد. این جنبش خود را جایگزینی برای راست سنتی (محافظه‌کاری، جمهوری‌خواهان) می‌داند. ویژگی‌های آن:
\begin{itemize}
    \item ناسیونالیسم سفید یا هویت‌گرایی سفید
    \item ضد لیبرالیسم و ضد «صحت سیاسی»
    \item فرهنگ اینترنتی: میم، ترول، طنز تاریک
    \item ضد فمینیسم و ضد چندفرهنگی
\end{itemize}
\end{tcolorbox}

\subsection{تایم‌لاین ظهور آلت‌رایت}

\begin{center}
\begin{tikzpicture}[
    scale=0.9,
    every node/.style={font=\small}
]
% Timeline base
\draw[line width=2pt, color=rightblue] (0,0) -- (15,0);

% Year markers
\foreach \x/\year in {0/2008, 3/2010, 6/2012, 9/2014, 12/2016, 15/2018} {
    \draw[line width=1.5pt, color=rightblue] (\x,-0.2) -- (\x,0.2);
    \node[below] at (\x,-0.3) {\year};
}

% Events
\node[above, align=center, fill=goldaccent!20, rounded corners, inner sep=3pt] at (1,0.5) {
    \begin{tabular}{c}
    {\footnotesize\rl{ریچارد اسپنسر}}\\
    {\footnotesize\rl{واژهٔ آلت‌رایت}}
    \end{tabular}
};

\node[above, align=center, fill=rightblue!20, rounded corners, inner sep=3pt] at (5,1.5) {
    \begin{tabular}{c}
    {\footnotesize\rl{گیمرگیت}}\\
    {\footnotesize\rl{جنگ فرهنگی آنلاین}}\\
    {\scriptsize 2014}
    \end{tabular}
};

\node[above, align=center, fill=leftred!20, rounded corners, inner sep=3pt] at (9,0.5) {
    \begin{tabular}{c}
    {\footnotesize\rl{نامزدی ترامپ}}\\
    {\footnotesize\rl{رشد آلت‌رایت}}
    \end{tabular}
};

\node[above, align=center, fill=centergreen!20, rounded corners, inner sep=3pt] at (12,1.5) {
    \begin{tabular}{c}
    {\footnotesize\rl{پیروزی ترامپ}}\\
    {\footnotesize\rl{اوج آلت‌رایت}}
    \end{tabular}
};

\node[above, align=center, fill=darkgray!20, rounded corners, inner sep=3pt] at (14,0.5) {
    \begin{tabular}{c}
    {\footnotesize\rl{شارلوتزویل}}\\
    {\footnotesize\rl{بحران و افول}}
    \end{tabular}
};

% Title
\node[above, font=\bfseries\large] at (7.5,3) {\rl{تایم‌لاین: ظهور و افول آلت‌رایت}};

\end{tikzpicture}
\end{center}

\subsection{فرهنگ اینترنتی و جنگ میم‌ها}


آلت‌رایت پدیده‌ای ذاتاً اینترنتی است. این جنبش در فضاهایی مانند ۴chan، ردیت، و توییتر شکل گرفت و از ابزارهای فرهنگ اینترنتی استفاده کرد:

\textbf{میم‌ها:} تصاویر و ویدئوهای طنزآمیز که پیام‌های سیاسی را منتقل می‌کنند. «پِپِ قورباغه» نماد غیررسمی آلت‌رایت شد.

\textbf{ترولینگ:} آزار و اذیت آنلاین مخالفان. آلت‌رایتی‌ها این کار را «مزاح» و «آزادی بیان» می‌نامند.

\textbf{زبان رمزی:} استفاده از کدها و اشارات برای فرار از سانسور. مثلاً «۱۴۸۸» اشاره به شعارهای نازی دارد.

\textbf{Irony و Post-irony:} آلت‌رایتی‌ها اغلب ادعا می‌کنند که «فقط شوخی می‌کنند». این ابهام عامدانه، تشخیص باورهای واقعی را دشوار می‌کند.

این استراتژی به آلت‌رایت امکان داد ایده‌هایی را که در گفتمان عمومی غیرقابل قبول بودند، به شکل «طنز» و «فرهنگ جوانان» معرفی کند.


\subsection{ارتباط با ترامپیسم}


آلت‌رایت خود را حامی دونالد ترامپ می‌دانست، هرچند ترامپ رسماً از آنها فاصله گرفت.

\textbf{نقاط اشتراک:}
\begin{itemize}
    \item ضد «نخبگان» و «استابلیشمنت»
    \item ضد مهاجرت غیرقانونی
    \item ضد «صحت سیاسی»
    \item ضد رسانه‌های جریان اصلی
    \item ناسیونالیسم اقتصادی
\end{itemize}

\textbf{استیو بنن}، مدیر کمپین ترامپ و سپس مشاور ارشد او، پل ارتباطی میان ترامپ و آلت‌رایت بود. بنن سردبیر سایت «برایتبارت» بود که آن را «پلتفرم آلت‌رایت» نامید.

پس از حادثهٔ شارلوتزویل در ۲۰۱۷ (که در آن یک ناسیونالیست سفید با ماشین به معترضان زد و یک نفر کشته شد)، آلت‌رایت به بحران فرورفت. بسیاری از پلتفرم‌ها آنها را مسدود کردند و جنبش تکه‌تکه شد.


%═══════════════════════════════════════════════════════════════
\section{پوپولیسم راست}
%═══════════════════════════════════════════════════════════════

\begin{tcolorbox}[
    enhanced,
    colback=rightblue!8,
    colframe=rightblue,
    boxrule=1.5pt,
    arc=5pt,
    fonttitle=\bfseries,
    title={\rl{🔵 تعریف کلیدی: پوپولیسم}}
]
\textbf{پوپولیسم} سبک سیاسی است که جامعه را به دو گروه متخاصم تقسیم می‌کند: «مردم» خوب و «نخبگان» فاسد. پوپولیست مدعی است که صدای واقعی مردم است و نخبگان منافع مردم را نادیده می‌گیرند.

پوپولیسم می‌تواند چپ یا راست باشد:
\begin{itemize}
    \item \textbf{پوپولیسم چپ:} مردم عادی علیه نخبگان اقتصادی (سرمایه‌داران)
    \item \textbf{پوپولیسم راست:} مردم «واقعی» علیه نخبگان فرهنگی و سیاسی + مهاجران
\end{itemize}
\end{tcolorbox}

\subsection{ترامپ، برگزیت، اوربان}


در دههٔ ۲۰۱۰، پوپولیسم راست به پیروزی‌های بزرگی دست یافت:

\textbf{ترامپ (۲۰۱۶):}
\begin{itemize}
    \item شعار: «آمریکا را دوباره بزرگ کن»
    \item ضد مهاجرت: «دیوار» در مرز مکزیک
    \item ضد نخبگان: «باتلاق را خشک کن»
    \item ضد جهانی‌شدن: «آمریکا اول»
\end{itemize}

\textbf{برگزیت (۲۰۱۶):}
\begin{itemize}
    \item شعار: «کنترل را بازپس بگیر»
    \item ضد بروکسل و بوروکراسی
    \item ضد مهاجرت آزاد
    \item ناسیونالیسم بریتانیایی
\end{itemize}

\textbf{اوربان (از ۲۰۱۰):}
\begin{itemize}
    \item «دموکراسی غیرلیبرال»
    \item ضد مهاجرت: حصار در مرز
    \item ضد سوروس و «نخبگان جهانی»
    \item ارزش‌های مسیحی و خانواده
\end{itemize}


\subsection{پوپولیسم: تهدید یا تصحیح؟}

\begin{tcolorbox}[
    enhanced,
    colback=leftred!8,
    colframe=leftred,
    boxrule=1.5pt,
    arc=5pt,
    fonttitle=\bfseries,
    title={\rl{🔴 مناظره: پوپولیسم}}
]

\textbf{پوپولیسم تهدید است:}
\begin{itemize}
    \item دموکراسی را تضعیف می‌کند
    \item نهادها را زیر سؤال می‌برد
    \item جامعه را قطبی می‌کند
    \item اقلیت‌ها را هدف می‌گیرد
    \item راه‌حل‌های ساده‌انگارانه می‌دهد
    \item به اقتدارگرایی می‌انجامد
\end{itemize}



\textbf{پوپولیسم تصحیح است:}
\begin{itemize}
    \item صدای فراموش‌شدگان را می‌رساند
    \item نخبگان را پاسخگو می‌کند
    \item مشارکت سیاسی را افزایش می‌دهد
    \item نشانهٔ بحران دموکراسی است
    \item مشکلات واقعی را نشان می‌دهد
    \item می‌تواند اصلاحات بیاورد
\end{itemize}

\end{tcolorbox}

%═══════════════════════════════════════════════════════════════
\section{محافظه‌کاری ملی}
%═══════════════════════════════════════════════════════════════


«محافظه‌کاری ملی» (National Conservatism) جریانی است که از ۲۰۱۹ به طور رسمی سازماندهی شده. این جریان تلاش می‌کند بین محافظه‌کاری سنتی، ناسیونالیسم، و نقد نئولیبرالیسم پل بزند.

\textbf{اصول محافظه‌کاری ملی:}
\begin{itemize}
    \item \textbf{دولت-ملت:} واحد سیاسی بنیادین است
    \item \textbf{نقد جهانی‌شدن:} از راست! نه به خاطر نابرابری، بلکه به خاطر تضعیف ملت
    \item \textbf{حمایت‌گرایی:} تجارت آزاد مطلق مضر است
    \item \textbf{محافظه‌کاری اجتماعی:} خانواده، دین، سنت
    \item \textbf{نقد لیبرتارینیسم:} بازار همه چیز نیست
\end{itemize}

این جریان با لیبرتارینیسم و نئوکنسرواتیسم تفاوت دارد:


\begin{table}[H]
\centering
\begin{tabular}{|>{\columncolor{rightblue!10}}r|p{3cm}|p{3cm}|p{3cm}|}
\hline
\rowcolor{rightblue!30}
\textbf{موضوع} & \textbf{محافظه‌کاری ملی} & \textbf{نئوکنسرواتیسم} & \textbf{لیبرتارینیسم} \\
\hline
\textbf{تجارت} & حمایت‌گرایی & تجارت آزاد & تجارت آزاد \\
\hline
\textbf{مهاجرت} & محدود & نسبتاً باز & آزاد یا محدود \\
\hline
\textbf{سیاست خارجی} & غیرمداخله‌جویانه & مداخله‌جویانه & غیرمداخله‌جویانه \\
\hline
\textbf{دولت} & دولت قوی ملی & دولت قوی جهانی & دولت حداقلی \\
\hline
\textbf{فرهنگ} & ملی و سنتی & غربی و دموکراتیک & فردگرا و خنثی \\
\hline
\end{tabular}
\caption{مقایسهٔ جریان‌های راست معاصر}
\end{table}

%═══════════════════════════════════════════════════════════════
\section{نئوری‌اکشنیسم و نقد دموکراسی}
%═══════════════════════════════════════════════════════════════


در حاشیهٔ راست نو، جریانی رادیکال‌تر وجود دارد که دموکراسی را به طور بنیادین زیر سؤال می‌برد.

\textbf{کرتیس یاروین} (متولد ۱۹۷۳)، با نام مستعار «مِنسیوس مولدباگ»، بلاگری است که ایدهٔ \textbf{نئوری‌اکشنیسم} (Neoreaction) را مطرح کرد.

\textbf{ایده‌های یاروین:}
\begin{itemize}
    \item دموکراسی ذاتاً ناکارآمد است
    \item «کلیسای تاریک» (نخبگان لیبرال) بر جامعه حکومت می‌کند
    \item بازگشت به شکلی از سلطنت یا «CEO-کراسی»
    \item کشور باید مانند شرکت اداره شود
    \item «خروج» به جای «رأی»: اگر نمی‌پسندی، برو
\end{itemize}

یاروین در سیلیکون ولی نفوذ دارد و پیتر تیل از حامیان اوست. این ایده‌ها، هرچند حاشیه‌ای، نشان‌دهندهٔ گسست برخی در راست از دموکراسی لیبرال است.


%═══════════════════════════════════════════════════════════════
\section{صنعت فکری راست نو}
%═══════════════════════════════════════════════════════════════

\subsection{رسانه‌های جایگزین}


راست نو شبکه‌ای از رسانه‌ها و اتاق‌های فکر ساخته که با «رسانه‌های جریان اصلی» رقابت می‌کنند:

\textbf{رسانه‌ها:}
\begin{itemize}
    \item \textbf{فاکس نیوز:} راست محافظه‌کار سنتی
    \item \textbf{برایتبارت:} راست پوپولیست/آلت‌رایت
    \item \textbf{دیلی وایر:} محافظه‌کاری جوان
    \item \textbf{پادکست‌ها و یوتیوب:} جو روگان، بن شاپیرو
\end{itemize}

\textbf{اتاق‌های فکر:}
\begin{itemize}
    \item هریتیج فاوندیشن (سنتی)
    \item کلرمونت اینستیتوت (ترامپیست)
    \item هازونی و ادموند برک فاوندیشن (محافظه‌کاری ملی)
\end{itemize}


\subsection{اینفلوئنسرهای راست}


\textbf{جردن پیترسون} (متولد ۱۹۶۲)، روانشناس کانادایی، با مخالفت با قانون C-16 کانادا (حمایت از تراجنسیتی‌ها) به شهرت رسید. او:
\begin{itemize}
    \item از «مردانگی سنتی» دفاع می‌کند
    \item «مارکسیسم فرهنگی» را نقد می‌کند
    \item کتاب \textit{۱۲ قانون زندگی} پرفروش شد
    \item جوانان مرد را جذب کرده
\end{itemize}

\textbf{بن شاپیرو} (متولد ۱۹۸۴)، مفسر سیاسی آمریکایی:
\begin{itemize}
    \item محافظه‌کار سنتی، نه آلت‌رایت
    \item مناظره‌های سریع و تهاجمی
    \item «حقایق به احساسات اهمیت نمی‌دهند»
    \item پادکست پرشنونده
\end{itemize}

آیا این افراد «راست افراطی» هستند؟ موضوع بحث است. آنها خود را محافظه‌کار سنتی می‌دانند، اما منتقدان می‌گویند «دروازهٔ» ورود به ایده‌های افراطی‌ترند.


%═══════════════════════════════════════════════════════════════
\section*{خلاصهٔ فصل}
%═══════════════════════════════════════════════════════════════

\begin{tcolorbox}[
    enhanced,
    colback=lightbg,
    colframe=darkgray,
    boxrule=2pt,
    arc=8pt,
    fonttitle=\bfseries\large,
    title={\rl{📝 خلاصهٔ فصل سیزدهم}}
]
\begin{itemize}
    \item \textbf{راست نو اروپایی} (آلن دوبنوا) استراتژی «متاپولیتیک» را اتخاذ کرد: تغییر فرهنگ پیش از سیاست. این جریان ضد لیبرالیسم، ضد مسیحیت، و ضد آمریکاست.
    
    \item \textbf{آلت‌رایت آمریکایی} جنبشی اینترنتی بود که از میم‌ها و طنز تاریک استفاده می‌کرد. این جنبش پس از شارلوتزویل (۲۰۱۷) تکه‌تکه شد.
    
    \item \textbf{پوپولیسم راست} (ترامپ، برگزیت، اوربان) «مردم» را علیه «نخبگان» و مهاجران بسیج می‌کند.
    
    \item \textbf{محافظه‌کاری ملی} تلاش می‌کند ناسیونالیسم، محافظه‌کاری اجتماعی، و نقد نئولیبرالیسم را ترکیب کند.
    
    \item راست نو شبکهٔ گسترده‌ای از \textbf{رسانه‌ها و اینفلوئنسرها} دارد که با رسانه‌های جریان اصلی رقابت می‌کنند.
    
    \item پرسش کلیدی: آیا این جریان‌ها «فاشیسم جدید» هستند یا پاسخی مشروع به مشکلات واقعی؟
\end{itemize}
\end{tcolorbox}

%═══════════════════════════════════════════════════════════════
\section*{پرسش‌های تأملی}
%═══════════════════════════════════════════════════════════════

\begin{tcolorbox}[
    enhanced,
    colback=centergreen!8,
    colframe=centergreen,
    boxrule=1.5pt,
    arc=5pt,
    fonttitle=\bfseries,
    title={\rl{❓ پرسش‌هایی برای تأمل و بحث}}
]
\begin{enumerate}
    \item آیا استراتژی «متاپولیتیک» راست نو موفق بوده است؟ آیا فرهنگ جوامع غربی به راست چرخیده؟
    
    \item چرا آلت‌رایت از فرهنگ اینترنتی و میم استفاده کرد؟ آیا این استراتژی قابل تقلید است؟
    
    \item تفاوت «پوپولیسم» و «دموکراسی» چیست؟ آیا پوپولیست‌ها «صدای مردم» هستند یا تهدیدی برای دموکراسی؟
    
    \item آیا جردن پیترسون و بن شاپیرو «دروازهٔ» راست افراطی هستند، یا محافظه‌کاران مشروع؟
    
    \item چرا راست نو از مسیحیت فاصله می‌گیرد، در حالی که پوپولیست‌ها (اوربان) از «ارزش‌های مسیحی» دفاع می‌کنند؟
    
    \item آیا نقد «صحت سیاسی» موجه است، یا پوششی برای نژادپرستی و تبعیض؟
\end{enumerate}
\end{tcolorbox}

\vspace{10mm}
\begin{center}
\rule{0.6\textwidth}{1pt}

\vspace{3mm}
\textit{«وقتی نمی‌دانی دشمنت کیست، از همه می‌ترسی. وقتی می‌دانی دشمنت کیست، فقط از او می‌ترسی.»}

\vspace{2mm}
--- کارل اشمیت

\vspace{3mm}
\rule{0.6\textwidth}{1pt}
\end{center}
